\section{Models and Methods}
% Talk about overview of the model. What languages, packages, what it does overall, why it does it, briefly how it does it, what does it enable in the context of the thesis objectives.
% Novelty of modelling framework lies in the plume aggregation approach and the application of a photochemical trajectory model to simulate chemical evolution of aircraft plumes on a high resolution grid.
% This opens up new possibilities to investigate the cumulative physical and chemical effects that occur locally in aircraft plumes.

\subsection{Trajectory Model (traj\_gen)}
%Generates aircraft trajectories from flight plan data.
%Samples and interpolates waypoints to create smooth trajectories.
%Waypoints serve as emission release points for dispersion modeling.

\subsection{Meteorology Model (gen\_met)}
%Integrates real-world meteorological data (wind, temperature, humidity).
%Provides boundary conditions for plume dispersion and advection.
%Captures wind speed, temperature, and turbulence effects on plume transport.

\subsection{Background Chemistry Model (gen\_bg\_chem)}
%Initializes baseline atmospheric chemistry (NOx, ozone, VOCs).
%Sets the starting conditions for how emissions perturb the atmosphere.
%Critical for evaluating NOx saturation and ozone production.

\subsection{Aircraft Performance Model (ac\_perf)}
%Computes fuel burn and emissions based on flight conditions (altitude, thrust).
%Uses emission indices (EI) to calculate mass of NOx, CO2, and water vapor.
%Provides emission masses at each waypoint for further modeling.

\subsection{Emissions Model (emissions)}
%Quantifies the mass of species released during flight.
%Distributes emissions spatially and temporally along the trajectory.
%Provides key inputs for plume dispersion simulations.

\subsection{Dispersion/Advection Model (sim\_plumes)}
%Models the spread of emissions using a Gaussian plume framework.

\subsection{Plume-to-Grid Model (plume\_to\_grid)}
%The PL dataset generated previously provides 1D vector data on the temporal evolution of plume waypoints, regarding position, heading, plume cross-sectional dimensions and emissions masses. 
%To simulate the regional-scale chemical and physical processes which occur, and hence capture the cumulative effects due to plume intersections, it is necessary to aggregate the vector data to a common grid. 
%The Plume-to-Grid Model uses a novel conservative remapping method, which remaps Gaussian plumes onto a 2D Eulerian latitude-longitude grid. 



%Divides the plume into cross-sectional slices at each waypoint.
%Simulates how species such as NOx disperse spatially.
\subsubsection{Instantaneous Dispersion}
%Calculates species dispersion at each timestep and plume segment.
%Determines longitude and latitude of plume slice points.

\subsubsection{Gaussian Plume Aggregation}
% estimate the mass of emissions species which is aggregated to the 2-D Eulerian grid. 
% For each timestep, emissions species, and plume segment, apply the aggregation scheme.
% Calculate plume segment edge coordinates.
% Create a spatial bounding box around each plume segment.
%  Create a segment grid that contains all cells in segment bounding box.
%  Gaussian is discretised 

%Slice Definition:
%Discretizes the Gaussian plume into slices using quantiles.
%Defines plume width based on Gaussian distribution parameters.
%Creates bounding boxes for plume slices to limit grid cell search space.

%Slice Edge Calculation:
%Computes longitude and latitude of slice edges based on plume orientation.
%Provides precise positioning of plume on the grid.

%Slice Polygon Creation: 
%Create Shapely Polygon of each slice, 

%Slice Grid Calculation:
%Intersects plume slices with Eulerian grid cells.
%Uses geometric operations (Shapely Intersection) to calculate overlap area.
%Assigns weights based on the proportion of plume mass in each grid cell.

%Segment Grid Aggregation:
%Aggregates all slice grids to form the segment grid for each plume segment.
%Represents total mass distribution of the segment on the grid.

%Global Grid Update:
%Adds the segment grid to the global Eulerian grid.
%Repeats for all plume segments, timesteps, and species.

\subsection{Contrail Model (run\_cc)}
%Simulates formation and radiative impact of contrails.
%Accounts for contrail interactions with radiation (e.g., radiative shielding).
%Captures effects of overlapping contrails on radiation and climate.

\subsection{Chemistry Model (run\_boxm)}
%Uses a photochemical box model with 219 species and 606 reactions.
%Simulates NOx chemistry, ozone production, and hydroxyl radical reactions.
%Focuses on NOx saturation and reduced ozone production due to high NOx levels.
%Dynamically coupled with the dispersion model to evolve species over time.

\subsection{Validation}

%Mass Conservation:
%Ensures total emission mass is conserved in the grid remapping process.


\subsubsection{Mass Conversation}

%Computational Optimization:
%Uses bounding boxes to limit grid cell search.
%Parallelizes calculations to enhance performance.

\subsubsection{Box Model Outputs}


